%------------------------------------------------------------------------------
% \PART 1
%------------------------------------------------------------------------------
\chapter[Обзор существующих нейронных сетей для распознавания рукописных цифр]
{ОБЗОР СУЩЕСТВУЮЩИХ НЕЙРОННЫХ СЕТЕЙ ДЛЯ РАСПОЗНАВАНИЯ РУКОПИСНЫХ ЦИФР}

%------------------------------------------------------------------------------
% \PART 1.1 Text
%------------------------------------------------------------------------------
\section{Нейронные сети}\par
\hspace*{12.5 mm}Нейронная сеть представляет из себя совокупность нейронов,
соединенных друг с другом определенным образом. Нейрон представляет из себя 
элемент, который вычисляет выходной сигнал из совокупности входных сигналов.
Структура нейрона представлена на рисунке~\ref{fig:neiron}.

\insertfigure{neiron}{pic/neiron.png}{Общее представление нейрона}{9.0cm}

Нейроны выполняют следующую последовательность действий:

– приём входных сигналов;

– комбинирование входных сигналов;

– вычисление выходных сигналов;

– передача выходных сигналов следующим нейронам.

Между собой нейроны могут быть соединены абсолютно по-разному, это
определяется структурой конкретной сети. Искусственные нейронные сети 
можно рассматривать как направленный граф со взвешенными связями, в котором 
искусственные нейроны являются узлами. По архитектуре связей они могут быть 
сгруппированы в два класса: сети прямого распространения, в которых графы не 
имеют петель, и рекуррентные сети (RNN), или сети с обратными связями.

%------------------------------------------------------------------------------
% \PART 1.2 Text
%------------------------------------------------------------------------------
\section{Нейронные сети для распознавания рукописных цифр}\par
\hspace*{12.5 mm}Распознавание рукописных цифр является одной из классических 
задач машинного обучения, для решения которой применяется широкий спектр 
нейронных сетей. В зависимости от сложности задачи, требований к точности, 
скорости обработки и аппаратных ограничений используются различные архитектуры, 
включая полносвязные сети (FCNN), сверточные нейронные сети (CNN) и более 
сложные гибридные модели.

Одним из первых методов распознавания рукописных цифр стали многослойные 
перцептроны (MLP), относящиеся к классу полносвязных нейронных сетей. 
Такие сети состоят из входного слоя, нескольких скрытых слоев и выходного слоя, 
где каждый нейрон соединен со всеми нейронами предыдущего и последующего слоев. 
Они способны успешно решать задачу распознавания, но 
требуют большого количества параметров и вычислительных ресурсов, что делает 
их менее эффективными.

Более эффективными 
оказались сверточные нейронные сети. Эти сети включают сверточные слои, 
которые извлекают характерные признаки из изображения. CNN значительно 
превосходят полносвязные сети в точности и скорости работы, поскольку 
используют пространственные зависимости в данных и требуют меньше параметров.

Хотя рекуррентные нейронные сети и их усовершенствованные версии, такие 
как Long Short-Term Memory Network (LSTM) и Gated Recurrent Unit Network (GRU),
чаще применяются для обработки последовательных данных, они могут 
использоваться и для распознавания рукописных цифр. Они 
эффективны в случаях, когда рукописные символы анализируются в контексте 
строки.

Для повышения точности и эффективности могут применяться гибридные подходы, 
комбинирующие преимущества разных типов нейронных сетей. Например, модели, 
совмещающие CNN и LSTM, успешно применяются для распознавания рукописного 
текста, где CNN используется для выделения признаков, а LSTM – для анализа 
последовательностей.
Рассмотрим подробнее полносвязные и сверточные нейронные сети.

%------------------------------------------------------------------------------
% \PART 1.3 Text
%------------------------------------------------------------------------------
\section{Полносвязные нейронные сети}\par
\hspace*{12.5 mm}Полносвязная нейронная сеть – 
это базовая архитектура искусственных нейронных сетей, в которой каждый нейрон 
одного слоя соединен со всеми нейронами следующего слоя. Такая архитектура
чаще всего применяется в качестве классификатора после сверточных или 
рекуррентных слоев. Количество нейронов на каждом слое может отличаться от 
соседних слоев.

Общая структура полносвязных нейронных сетей показана на 
рисунке~\ref{fig:fcnn}.

\insertfigure{fcnn}{pic/fcnn.png}{Общая структура полносвязных нейронных сетей}{12.0cm}

Математически данную сеть можно описать формулой для вычисления выхода нейрона 
в полносвязном слое
\begin{equation}
    {h} = f({W} {x} + {b}),
\end{equation}

\noindentгде   ${x}$ – входной вектор;
 ${W}$ – матрица весов;
 ${b}$ – вектор смещений;
 $f(\cdot)$ – функция активации;
 ${h}$ – выходной вектор нейронов.

%------------------------------------------------------------------------------
% \PART 1.4 Text
%------------------------------------------------------------------------------
\section{Сверточные нейронные сети}
\hspace*{12.5 mm}
Сверточные нейронные сети (Convolutional Neural Networks) представляют собой 
архитектуру глубокого обучения, предназначенную для обработки данных обладающих 
сетчатой топологией, таких как изображения. Их основное преимущество 
заключается в способности эффективно извлекать пространственные и иерархические 
особенности входных данных с помощью операций свертки.

Принцип работы сверточных нейронных сетей показан на рисунке~\ref{fig:cnn}\cite{CNN}.

\insertfigure{cnn}{pic/cnn.png}{Принцип работы сверточных нейронных сетей}{10cm}

Основная операция в сверточных нейронных сетях – это свертка, 
которая заключается в пошаговом применении фильтра (ядра 
свёртки) к локальным областям входного изображения:

\begin{equation}
    S(i, j) = \sum_{m} \sum_{n} X(i+m, j+n) \cdot K(m, n),
\end{equation}

\noindentгде  \( X(i, j) \) – входное изображение;
\( K(m, n) \) – ядро свертки;
\( S(i, j) \) – выходное изображение после свертки.

Сверточный слой применяется к входному изображению, используя несколько 
фильтров, каждый из которых извлекает определенные характеристики. Как правило,
чем глубже слой, тем более абстрактные признаки он способен захватывать.
